\documentclass[11pt]{article}
\usepackage{amsmath,amssymb,amsthm}
\usepackage[margin=1.1in]{geometry}
\newtheorem{lemma}{Lemma}
\newtheorem{proposition}{Proposition}
\newtheorem{theorem}{Theorem}
\theoremstyle{remark}
\newtheorem{remark}{Remark}
\newcommand{\fl}[1]{\left\lfloor #1 \right\rfloor}

\title{Prime counts and prime sums from multiplication tables:\\ a self-similar formula system}
\author{Carl Gribble\thanks{Draft v0.1, companion to \emph{Certified exact sums of primes over dyadic
intervals via a ladder of generalized factorials}~\cite{Companion} (same author). Developed in collaboration with Claude
(Anthropic); all numerical claims were machine-verified as described in Section~6. Responsibility for
the content rests with the author. A full disclosure of AI use appears at the end of the paper.}}
\date{August 2026 --- draft}

\begin{document}
\maketitle

\begin{abstract}
A multiplication table whose entries are all at least $2$ manufactures exactly the composite numbers;
the primes are precisely what the table cannot make. Taking this observation seriously as a computational
program --- ``sum of the integers, minus the de-duplicated output of the table, equals the sum of the
primes'' --- forces a de-duplication problem, since a composite $m$ appears $d(m)-2$ times in the table,
and $2^{\Omega(m)}$-fold growth in higher-dimensional tables. Its exact resolution is the
alternating-harmonic combination of $k$-dimensional multiplication tables, which is Linnik's identity.
We develop this into a three-line inductive formula family $S_j$ computing $\sum_{p \le y} p^{\,j}$ for
every weight $j \ge 0$ --- counts at $j=0$, sums at $j=1$, power sums beyond --- over arbitrary
intervals, with a self-similar ``dust'' recursion that terminates in \emph{empty} tables, so that no
prime, primality test, or sieve appears anywhere in the system, at any depth. We prove the interval
structure: sums live on hyperbola shells, folded coordinates are root-bounded, evaluation requires only
the finite quotient--root lattice of the endpoint, and dimension $k$ is \emph{born} at $2^k$ with a
single newborn. We prove a duplicated-mass theorem: the $2$-dimensional table cancels identically out
of its own duplication ledger, which is settled by the diagonal and the tables of dimension $\ge 3$
alone. A reference implementation passes forty exact interval tests for counts and sums, and at the
billion scale reproduces $\pi(10^9) = 50{,}847{,}534$ together with the primality and index of
$999{,}999{,}937$ in minutes, touching $63{,}719$ lattice points out of $10^9$ integers. A curious
observable falls out: stepping the endpoint from $y$ to $y+1$ enlarges the evaluation lattice exactly at
the divisors of $y+1$, so a prime step adds a single point --- primality is visible as minimal lattice
growth. The core identity is Linnik's, the resulting algorithms sit in the Legendre--Meissel family, and the
computational route through Linnik's identity was pioneered by McKenzie; the contribution is the
geometric derivation, the duplicated-mass and layer-birth theorems, the interval taxonomy, and the
verification discipline.
\end{abstract}

\medskip
\noindent\textbf{Keywords:} prime counting function; prime power sums; Linnik's identity;
multiplication tables; divisor function; Legendre--Meissel methods; machine-verified computation.

\smallskip
\noindent\textbf{MSC 2020:} primary 11N05; secondary 11A25, 11Y35, 11Y16.

\section{Introduction}

Let $U_k(N)$ denote the $k$-dimensional multiplication table restricted to entries at least $2$ and
products at most $N$: the set of tuples $(x_1, \dots, x_k)$ with every $x_i \ge 2$ and
$x_1 x_2 \cdots x_k \le N$. The starting observation is elementary and, we believe, a good pedagogical
door into multiplicative number theory: the values manufactured by $U_2$ are exactly the composite
numbers, and the primes are exactly the integers the table cannot make. An arithmetic amateur's very
natural program --- \emph{sum the integers up to $N$, subtract what the table makes, obtain the sum of
the primes} --- is therefore structurally sound, and fails only on multiplicity: the table ships
duplicates, in quantities governed by the divisor function and, more fundamentally, by the number of
prime factors $\Omega(m)$, which we call the \emph{dimension} of $m$ throughout.

This paper carries the program to completion and then generalizes it. The exact de-duplication weights
turn out to be the alternating-harmonic sequence $1, -\tfrac12, +\tfrac13, \dots$ applied to the tower
$U_1, U_2, U_3, \dots$; per integer this is Linnik's identity \cite{Linnik,FI}, the combinatorial
expansion of $\log \zeta$. Around it we build a uniform, weighted, self-similar formula system with the
following properties, each verified computationally at scales up to $10^9$: it computes
$\sum_{A < p \le B} p^{\,j}$ exactly for every $j \ge 0$ and every interval; it consumes no prime data at
any point, its recursion bottoming out in empty tables; its interval structure is completely described
by hyperbola shells, root bounds, and a finite quotient--root lattice; and it exhibits a
\emph{layer-birth law} organizing the dyadic scale $2^k$ as the birthday of dimension $k$.

We state our claims and non-claims plainly. The per-integer identity is Linnik's; the $j = 0$
extraction line is Riemann's $\Pi$-to-$\pi$ inversion \cite{Riemann,Edwards} with $\Pi$ supplied by
table layers rather than an integral; and the resulting algorithms belong to the
Legendre--Meissel combinatorial family \cite{DR,Staple,Orlov,Lucy,McKenzie}, which contains asymptotically
faster members. What we believe is worth recording: the derivation of the identity from the
de-duplication ledger of a literal times table; the uniform treatment of counts, sums, and power sums as
one three-line induction with a weight dial; the duplicated-mass cancellation theorem
(Section~4); the interval taxonomy and layer-birth law (Section~5); the fully prime-free self-similar
dust (Section~3); and the lattice-growth observable at prime steps (Section~6). Of these, the computational core --- prime counts and
prime power sums from Linnik's identity, with no prime input --- is anticipated by McKenzie
\cite{McKenzie}, whose $x^{2/3}$-class algorithm predates and asymptotically improves on our
implementation; the remaining items we have not found presented in this form, though we regard that as
a statement about our searches.

\section{The table frame}

Write $D_k^*(n)$ for the number of ordered factorizations of $n$ into $k$ factors, each at least $2$;
thus $D_1^*(n) = 1$ for $n \ge 2$, $D_2^*(n) = d(n) - 2$, and $D_k^*(n) = 0$ unless $\Omega(n) \ge k$,
since dimension is additive and each factor carries dimension at least $1$. Two appearance laws anchor
intuition. In the (unrestricted) $k$-fold table the multiplicity of $n$ is
$d_k(n) = \prod_{p^a \| n} \binom{a + k - 1}{k - 1}$ --- the Euler coefficients of $\zeta(s)^k$ \cite{HW} --- so a
squarefree $n$ of dimension $\omega$ appears $k^{\omega}$ times: base the dimension of the space,
exponent the dimension of the number. In the restricted table, primes appear only in $U_1$: a prime
admits no factorization into two factors $\ge 2$. This is the survival mechanism of everything below.

The $x \leftrightarrow y$ fold of $U_2$ along its diagonal is Dirichlet's hyperbola device
\cite{Dirichlet}; the diagonal itself carries the perfect squares, with the closed form
$\sum_{k=2}^{s} k^2 = \tfrac{s(s+1)(2s+1)}{6} - 1$, $s = \fl{\sqrt N}$, and after folding, the residual
duplicates are exactly the values with several distinct two-part factorizations --- necessarily of
dimension $\ge 3$. (At $N = 100$: the folded interior square repeats only
$12, 16, 18, 20, 24, 30, 36, 40$, of dimensions $3,4,3,3,4,3,4,4$.)

\section{The de-duplication identity and the formula system}

\begin{lemma}[Per-integer ledger; Linnik]\label{lem:linnik}
For every integer $n \ge 2$,
\[
\sum_{k \ge 1} \frac{(-1)^{k-1}}{k}\, D_k^*(n) \;=\; \frac{\Lambda(n)}{\log n} \;=\;
\begin{cases} 1/j, & n = p^{\,j},\\ 0, & \text{otherwise,}\end{cases}
\]
the sum terminating at $k = \Omega(n)$.
\end{lemma}

The generating identity is $\log \zeta = \log(1 + F) = F - \tfrac{F^2}{2} + \cdots$ with
$F = \zeta - 1$; the ledger of each composite closes at its own dimension, and primes survive with
weight $1$ because they inhabit only $U_1$. An equivalent integer-weight system uses dimension parity,
$\mu(d) = (-1)^{\Omega(d)}$ on squarefree sieving patterns (Legendre's sieve); we work with the harmonic
weights throughout.

Fix a weight $j \ge 0$ and define the layer totals $T_k^{(j)}(y) = \sum_{n \le y} n^j D_k^*(n)$. The
whole system is three lines:
\begin{align}
\text{base:}\quad & T_1^{(j)}(y) \;=\; \textstyle\sum_{n=2}^{y} n^j \quad\text{(Faulhaber closed form)},\label{eq:base}\\
\text{step:}\quad & T_k^{(j)}(y) \;=\; \sum_{d=2}^{\fl{y/2^{k-1}}} d^{\,j}\; T_{k-1}^{(j)}\big(\fl{y/d}\big),\label{eq:step}\\
\text{extraction:}\quad & S_j(y) \;=\; \sum_{k=1}^{\fl{\log_2 y}} \frac{(-1)^{k-1}}{k}\, T_k^{(j)}(y)
\;-\; \sum_{a=2}^{\fl{\log_2 y}} \frac{1}{a}\; S_{a j}\big(\fl{y^{1/a}}\big).\label{eq:extract}
\end{align}

\begin{theorem}[Prime power sums, prime-free]\label{thm:main}
For every $j \ge 0$ and $y \ge 1$, $S_j(y) = \sum_{p \le y} p^{\,j}$; in particular $S_0 = \pi$ and
$S_1$ is the prime-sum function, and for any interval,
$\sum_{A < p \le B} p^{\,j} = S_j(B) - S_j(A)$. The recursion \eqref{eq:extract} is self-similar (the
dust calls the same $S$ at root scales with multiplied weight) and terminates without any prime input:
at $y \le 3$ all layers with $k \ge 2$ vanish because the tables are empty, and the formula returns
$\sum_{n \le y} n^j$ over $\{2, 3\} \cap [1, y]$ of its own accord.
\end{theorem}

The alternating combination in \eqref{eq:extract} is Riemann's prime-power counting function (weighted),
and the extraction line is his $\Pi$-to-$\pi$ inversion made recursive; the recursion depth is
$O(\log \log y)$. As a worked instance, at $y = 30$, $j = 0$: the layers are $29, 52, 32, 5$, giving
$29 - \tfrac{52}{2} + \tfrac{32}{3} - \tfrac{5}{4} = \tfrac{149}{12} = \Pi(30)$; the dust is
$\tfrac{\pi(5)}{2} + \tfrac{\pi(3)}{3} + \tfrac{\pi(2)}{4} = \tfrac{29}{12}$; and
$\pi(30) = \tfrac{149}{12} - \tfrac{29}{12} = 10$. All rungs are linked by Abel summation, so counts
rebuild sums and sums rebuild counts; and unit windows give membership,
$S_1(n) - S_1(n-1) \in \{0, n\}$, so the aggregate formulas determine the set of primes itself by
bisection.

\section{The duplicated-mass theorem}

Return to the amateur's program in its corrected cage: the folded region $2 \le x \le \fl{\sqrt N}$,
$x \le y$, $xy \le N$, which manufactures every composite $\le N$. Its raw sum is
$R(N) = \tfrac12\big(T_2^{(1)}(N) + Q(N)\big)$ with $Q$ the diagonal closed form above; the duplicated
mass is $D(N) = R(N) - \sum_{m \le N \text{ composite}} m$.

\begin{theorem}[The 2D table cancels from its own duplication bill]\label{thm:dup}
\[
D(N) \;=\; \frac{Q(N)}{2} \;+\; \sum_{k \ge 3} \frac{(-1)^{k-1}}{k}\, T_k^{(1)}(N)
\;-\!\!\sum_{p^{j} \le N,\ j \ge 2} \frac{p^{\,j}}{j}.
\]
\end{theorem}

The proof is three lines from Lemma~\ref{lem:linnik}: the composite indicator is
$\tfrac{D_2^*}{2} - \tfrac{D_3^*}{3} + \cdots$ up to prime-power dust, and upon subtracting, the
$T_2^{(1)}$ terms annihilate. The duplication of the two-dimensional table is settled entirely by the
squares and the tables of dimension three and higher --- in agreement with the hand ledger, where every
post-fold repeat offender had dimension $\ge 3$. Verified exactly: $D(100) = 4593$ from
$192 + \tfrac{21426}{3} - \tfrac{13441}{4} + \tfrac{4112}{5} - \tfrac{640}{6} - \tfrac{5789}{60}$, and
$D(1000) = 974{,}163$. The dust term is itself produced prime-free by Theorem~\ref{thm:main}.

\section{Interval structure}

\begin{proposition}[Where the intervals live]\label{prop:intervals}
(i) For a target interval $(A, B]$ of values, every layer is summed over the hyperbola shell
$A < x_1 \cdots x_k \le B$; the interval constrains the product, never the coordinates. (ii) In the
sorted fold, the coordinates are root-bounded: $x_1 \le B^{1/k}$ and inductively
$x_i \le \big(B/(x_1 \cdots x_{i-1})\big)^{1/(k-i+1)}$. (iii) The system
\eqref{eq:base}--\eqref{eq:extract} evaluates only at the closure of $\{B\}$ (and $\{A\}$) under
$y \mapsto \fl{y/d}$ and $y \mapsto \fl{y^{1/a}}$ --- a finite, self-similar lattice of
$O(\sqrt B)$ points (e.g.\ $642$ points for $B = 10^5$).
\end{proposition}

\begin{proposition}[Layer-birth law]\label{prop:birth}
$T_k^{(j)}(y) = 0$ for $y < 2^k$, and the contribution of layer $k$ to its own dyadic window
$(2^{k-1}, 2^k]$ is the single tuple $(2, 2, \dots, 2)$, of value $2^k$ and multiplicity one.
Consequently the extraction truncates at $k = \fl{\log_2 y}$, and the dyadic window of index $k$ is the
birth window of dimension $k$, containing exactly one newborn.
\end{proposition}

\section{Implementation and verification}

A reference implementation (\texttt{prime\_sum\_formula.py} in the code repository \cite{Code}, about
one hundred lines)
realizes \eqref{eq:base}--\eqref{eq:extract} with exact rational arithmetic, memoization over the
quotient--root lattice, and hyperbola block-grouping of the step \eqref{eq:step}, so each lattice node
costs $O(\sqrt y)$; observed total cost scales near $y^{3/4}$. Every returned value is asserted to have
denominator $1$ before use --- the harmonic fractions must annihilate, and in all runs they did.

\emph{Test suite.} Forty interval tests, all exact against an independent sieve: the twenty dyadic
windows to $2^{20}$ (sums); edge cases including the empty interval, the prime-free window
$(113, 126]$, a unit window, and a window containing $2^{10}$ to exercise the dust; eight random
intervals below $10^6$; five count tests (weight $j = 0$), including $\pi(10^6) = 78{,}498$; and the
spans $(10^6, 10^7]$ and $(1, 10^7]$, the latter returning $3{,}203{,}324{,}994{,}356$.

{\emergencystretch=2em
\emph{Billion-scale validation.} With published values as the oracle
($\pi(10^9) = 50{,}847{,}534$; $999{,}999{,}937$ the largest prime below $10^9$), the implementation
computed $\pi(999{,}999{,}936) = 50{,}847{,}533$ in $176$ s, $\pi(999{,}999{,}937) = 50{,}847{,}534$ in
$1$ s, and $\pi(10^9) = 50{,}847{,}534$ in $18$ s, touching $63{,}719$ lattice points in total: the
formula's verdicts --- that $999{,}999{,}937$ is prime, is prime number $50{,}847{,}534$, and is the
last prime below one billion --- agree with the record values, with no sieve and no primality test in
the computation.\par}

\begin{remark}[Primality as lattice growth]
Since $\fl{(y+1)/d} > \fl{y/d}$ exactly when $d \mid y+1$, the evaluation lattice of $y+1$ exceeds that
of $y$ only at divisors of $y+1$. A prime step therefore adds a single lattice point --- observed above
as the $176\,\mathrm{s} \to 1\,\mathrm{s}$ collapse and a lattice increment of exactly one. The
system's cost profile feels primality before its arithmetic announces it.
\end{remark}

\section{Related work, limitations, open doors}

Lemma~\ref{lem:linnik} is Linnik's identity \cite{Linnik}, standard in sieve theory \cite{FI} and
developed into $x^{2/3}$-class algorithms for $\pi(x)$ and prime power sums by McKenzie
\cite{McKenzie}; the $j=0$ extraction is Riemann's inversion \cite{Riemann,Edwards}; Legendre's sieve is the
dimension-parity twin; and the combinatorial prime-counting lineage
\cite{DR,Staple} and its prime-power-sum extensions \cite{Orlov,Lucy} contain algorithms of better
exponent ($x^{2/3}$-class) than our $x^{3/4}$-class implementation. The companion paper \cite{Companion} develops a
different, certified route to interval prime sums via Chebyshev-weighted generalized factorials. The
present system computes any functional of primes defined by \emph{individual} membership --- intervals
here; residue classes are an open door, since the ledger survives residue-restricted weights and the
Faulhaber corrections remain closed over progressions --- but no weighting can express \emph{joint}
conditions across several integers (twin primes, gaps): the per-integer collapse cannot see
correlations, which is the field's frontier, not this construction's defect. Finally, the system is a
formula \emph{family} of $\sqrt N$-length evaluations, not a fixed-length closed form; the obstruction
to further compression is the standard one, that the exact fluctuations carry the zeta spectrum.

\section*{Disclosure of AI use}

In accordance with publisher policies on the use of artificial intelligence, the author discloses the
following. This paper was developed in an extensive collaboration with Claude (Anthropic; Claude Fable
5, accessed through the claude.ai interface, 2026). The direction of the work --- the founding research
programme of exact moment systems for interval primes, of which this paper is the first --- together
with all decisions about scope, claims, and submission, are the author's. Within that direction, the
AI system contributed substantially to the mathematical development and proofs, wrote the verification
code, drafted and revised the manuscript text, and assisted with literature search; the reason for its
use was to enable a single non-specialist author to develop, check, and document the material to a
publishable standard. Every numerical and mathematical claim was subsequently machine-verified by the
test suite and billion-scale validation reported in Section~6, with the verification code openly
available in the cited repository \cite{Code}. The author has reviewed the entire manuscript and takes
full responsibility for its content, including the accuracy of the mathematics and of all references.

\begin{thebibliography}{99}
\bibitem{Linnik} Yu.~V. Linnik, \emph{The Dispersion Method in Binary Additive Problems}, Amer. Math.
Soc., Providence, 1963.
\bibitem{FI} J. Friedlander and H. Iwaniec, \emph{Opera de Cribro}, Amer. Math. Soc. Colloq. Publ.
57, 2010.
\bibitem{Riemann} B. Riemann, \"Uber die Anzahl der Primzahlen unter einer gegebenen Gr\"osse,
\emph{Monatsber. Berlin. Akad.} (1859).
\bibitem{Edwards} H.~M. Edwards, \emph{Riemann's Zeta Function}, Academic Press, 1974.
\bibitem{Dirichlet} P.~G.~L. Dirichlet, \"Uber die Bestimmung der mittleren Werthe in der
Zahlentheorie, \emph{Abh. K\"onigl. Preuss. Akad. Wiss.} (1849), 69--83.
\bibitem{DR} M. Del\'eglise and J. Rivat, Computing $\pi(x)$: the Meissel, Lehmer, Lagarias, Miller,
Odlyzko method, \emph{Math. Comp.} \textbf{65} (1996), 235--245.
\bibitem{Staple} D.~B. Staple, The combinatorial algorithm for computing $\pi(x)$, arXiv:1503.01839
(2015).
\bibitem{Orlov} A. Orlov, The combinatorial method to compute the sum of the powers of primes,
arXiv:2111.15545 (2021).
\bibitem{Lucy} Lucy\_Hedgehog, post in the Project Euler Problem 10 discussion thread (2013).
\bibitem{HW} G.~H. Hardy and E.~M. Wright, \emph{An Introduction to the Theory of Numbers}, 6th ed.,
Oxford Univ. Press, 2008.
\bibitem{McKenzie} N. McKenzie, Computing the prime counting function with Linnik's identity,
manuscript (2011), available at \texttt{https://www.primecounting.com}.
\bibitem{Companion} C. Gribble, Certified exact sums of primes over dyadic intervals via a ladder
of generalized factorials, companion manuscript, Zenodo (2026), \texttt{doi:10.5281/zenodo.21769105}.
\bibitem{Code} C. Gribble, Reference implementations accompanying this paper and \cite{Companion},
code repository (2026), \texttt{https://github.com/carlgribble-caa/prime-moments}.
\end{thebibliography}

\end{document}
